\documentclass[11pt]{article}

\usepackage[T1]{fontenc}
\usepackage[utf8]{inputenc}
\usepackage{lmodern}
\usepackage{microtype}
\usepackage[margin=1in]{geometry}
\usepackage{amsmath,amssymb,amsthm,mathtools}
\usepackage{enumitem}
\usepackage{xcolor}
\usepackage{graphicx}
\usepackage[hidelinks]{hyperref}
\usepackage[nameinlink,capitalize,noabbrev]{cleveref}

\setlist{nosep}
\allowdisplaybreaks

\newtheorem{theorem}{Theorem}[section]
\newtheorem{proposition}[theorem]{Proposition}
\newtheorem{lemma}[theorem]{Lemma}
\newtheorem{corollary}[theorem]{Corollary}
\newtheorem{question}[theorem]{Selection problem}
\theoremstyle{definition}
\newtheorem{definition}[theorem]{Definition}
\newtheorem{status}[theorem]{Status statement}
\theoremstyle{remark}
\newtheorem{remark}[theorem]{Remark}

\newcommand{\N}{\mathbb N}
\newcommand{\K}{\mathbb K}
\newcommand{\cP}{\mathcal P}
\newcommand{\I}{\mathcal I}
\newcommand{\J}{\mathcal J}
\newcommand{\CR}{\mathcal{CR}}
\newcommand{\Fin}{\mathrm{Fin}}
\newcommand{\Exh}{\mathrm{Exh}}
\newcommand{\Id}{\mathcal I_{\mathrm d}}
\newcommand{\SOT}{\mathrm{SOT}}
\newcommand{\supp}{\operatorname{supp}}
\newcommand{\Ideal}{\operatorname{Ideal}}
\newcommand{\spanop}{\operatorname{span}}
\newcommand{\restr}{\mathbin{\upharpoonright}}
\newcommand{\eps}{\varepsilon}
\newcommand{\one}{\mathbf 1}

\title{Critical Ideals of Ideal Schauder Systems:\\
A Verified Obstruction and Operator-Theoretic Reductions}
\author{Research note prepared from an automatic-solver packet}
\date{June 19, 2026}

\begin{document}
\maketitle

\begin{abstract}
Kwela and Swaczyna ask whether every analytic $P$-ideal is the critical ideal
of an ideal Schauder system in a Banach space.  This note gives the strongest
conclusions obtained in a renewed attempt at that problem.  First, it repairs
and proves the packet's unconditional-support theorem: every ideal defined by
convergence of positive weighted subseries along a $1$-unconditional basis is
non-pathological.  This is not a new characterization; it is a transparent
special case of the Banach-representability theorem of
Borodulin-Nadzieja--Farkas--Plebanek.

For a general ideal Schauder system, the note then rewrites the problem in
terms of the residual projections $R_n=I-S_n$.  Combining this reformulation
with a Katetov characterization of non-pathological analytic $P$-ideals gives
an exact statistical strong-operator selection criterion and an equivalent
regular-matrix averaging criterion.  Two usable sufficient conditions follow.
In particular, if every positive set contains a positive subset on which
$(\|R_n\|)$ is bounded, then the critical ideal is non-pathological.  Using the
filter-dependent uniform boundedness principle of De Bondt and Vernaeve, this
holds whenever the dual filter is rapid$^+$.  Finally, a theorem of Kadets and
Manskova yields a necessary reciprocal-norm divergence condition on every
positive set.

No complete proof or counterexample to the source question was located or
obtained.  The exact remaining obstruction is a statistical strong-operator
selection problem for an unbounded nested family of projections.
\end{abstract}

\section*{Source question screenshot}

The source question is Question 11.2 in Adam Kwela and Jaroslaw Swaczyna,
\emph{Zoo of ideal Schauder bases}, arXiv:2508.21235.

\begin{center}
\includegraphics[width=\linewidth]{figures/open_problem_crop.png}
\end{center}

\section*{Status and idea}

This packet is a strengthened partial result and reduction package, not a
complete answer to Question 11.2.  The original unconditional-support
obstruction remains valid after the repairs below, but the more useful output
is the operator-theoretic reformulation: for an analytic $P$-ideal critical
ideal, non-pathology is equivalent to a diffuse statistical strong-operator
selection principle for the residual projections $R_n=I-S_n$.  The packet also
records two concrete consequences: bounded positive thinning implies
non-pathology, and every positive set must satisfy a reciprocal-norm
divergence condition.

\section{The source question and the status of this note}

Let $X$ be a Banach space over $\K\in\{\mathbb R,\mathbb C\}$, and let
$\widehat a=(a_n,a_n^*)_{n\in\N}\subseteq X\times X^*$ be biorthogonal:
$a_n^*(a_m)=\delta_{nm}$.  Put
\[
 S_nx=\sum_{i=0}^n a_i^*(x)a_i,
 \qquad
 R_nx=x-S_nx.
\]
For an ideal $\I$ on $\N$, the system is an $\I$-Schauder system if
$R_nx\to0$ along the dual filter $\I^*$ for every $x\in X$.  The least ideal
with this property is its \emph{critical ideal}, denoted $\CR(X,\widehat a)$.
Equivalently,
\begin{equation}\label{eq:critical-definition}
 \CR(X,\widehat a)
 =\Ideal\Bigl(\Fin\cup
   \bigl\{B(x,\eps):x\in X,\ \eps>0\bigr\}\Bigr),
 \qquad
 B(x,\eps)=\{n:\|R_nx\|\geq\eps\}.
\end{equation}
Here $\Ideal(\mathcal A)$ denotes the ideal generated by the family
$\mathcal A$.

Kwela and Swaczyna proved that every non-pathological analytic $P$-ideal is a
critical ideal and asked the following question \cite[Question 11.2]{KS}:
\begin{quote}
Is every analytic $P$-ideal of the form $\CR(X,\widehat a)$ for some Banach
space $X$ and some biorthogonal system $\widehat a$ with continuous coordinate
functionals?
\end{quote}
The published version of their paper still contains the question.  The
literature search carried out for this note did not locate a subsequent paper
or preprint resolving it.  Thus the statements below are presented as partial
progress and reductions, not as a claimed solution of Question 11.2.

\section{Preliminaries on analytic $P$-ideals}

A function $\phi:\cP(\N)\to[0,\infty]$ is a \emph{submeasure} if
$\phi(\varnothing)=0$, $\phi(\{n\})<\infty$, it is monotone, and
$\phi(A\cup C)\leq\phi(A)+\phi(C)$.  It is lower semicontinuous if
\[
 \phi(A)=\sup_m\phi(A\cap m).
\]
Set
\[
 \Exh(\phi)=
 \{A\subseteq\N:\lim_{m\to\infty}\phi(A\setminus m)=0\}.
\]
Solecki's theorem says that the analytic $P$-ideals are exactly the ideals
$\Exh(\phi)$ for lower semicontinuous submeasures $\phi$ \cite{Solecki}.

A lower semicontinuous submeasure $\phi$ is \emph{non-pathological} if
\begin{equation}\label{eq:nonpath-submeasure}
 \phi(A)=\sup\{\mu(A):\mu\text{ is a finitely supported measure and }
                         \mu\leq\phi\}
\end{equation}
for every $A\subseteq\N$.  An analytic $P$-ideal is non-pathological if it is
$\Exh(\phi)$ for such a $\phi$.

We shall use the Katetov order.  For ideals $\I$ on $M$ and $\J$ on $N$, write
$\I\leq_K\J$ if there is $f:N\to M$ such that $f^{-1}[C]\in\J$ for every
$C\in\I$.  Let
\[
 \Id=\left\{A\subseteq\N:\frac{|A\cap m|}{m}\longrightarrow0\right\}
\]
be the asymptotic-density-zero ideal.  The following characterization is the
main descriptive-set-theoretic input used later.

\begin{theorem}[Hru\v{s}\'ak; see Filip\'ow--Tryba]\label{thm:hrusak}
For an analytic $P$-ideal $\I$, the following are equivalent:
\begin{enumerate}[label=(\roman*)]
 \item $\I$ is non-pathological;
 \item $\I\restr A\leq_K\Id$ for every $A\notin\I$.
\end{enumerate}
Moreover, an arbitrary ideal is an intersection of regular matrix summability
ideals if and only if it satisfies condition \emph{(ii)}.
\end{theorem}

The first assertion is recorded as Theorem 8.4 and the matrix assertion as
Theorem 9.1 in \cite{FilipowTryba}.

\section{The unconditional support theorem}

The automatic packet proposed the theorem below.  Its conclusion is correct,
but four details require care: the subadditivity split must be disjoint;
$\phi(F)$ must first be identified with the norm of the full positive sum; the
Cauchy argument should use ordered tails; and the positive norming functional
must be justified.  We include a complete proof.

\begin{theorem}[Unconditional support ideals are non-pathological]
\label{thm:unconditional-support}
Let $X$ have a $1$-unconditional Schauder basis $(u_n)_{n\in\N}$, and let
$c_n>0$.  Define
\[
 \J_{X,c}=\left\{A\subseteq\N:
                  \sum_{n\in A}c_nu_n\text{ converges in }X\right\}
\]
and
\begin{equation}\label{eq:phi-support}
 \phi(A)=\sup\left\{
   \left\|\sum_{n\in F}c_nu_n\right\|:F\in[A]^{<\omega}\right\}.
\end{equation}
Then $\phi$ is a non-pathological lower semicontinuous submeasure and
\[
 \J_{X,c}=\Exh(\phi).
\]
\end{theorem}

\begin{proof}
Monotonicity and $\phi(\varnothing)=0$ are immediate, and
$\phi(\{n\})=c_n\|u_n\|<\infty$.  If $F\subseteq A\cup C$ is finite, use the
disjoint decomposition
\[
 F=(F\cap A)\,\dot\cup\,(F\setminus A),
 \qquad F\setminus A\subseteq C.
\]
The triangle inequality gives
\[
 \left\|\sum_{n\in F}c_nu_n\right\|
 \leq \phi(A)+\phi(C).
\]
Taking the supremum over $F$ proves subadditivity.  Also
\[
 \phi(A)=\sup_m\phi(A\cap m),
\]
so $\phi$ is lower semicontinuous.

Suppose first that $\phi(A\setminus m)\to0$.  If $p>q\geq m$, then
\[
 \left\|\sum_{\substack{n\in A\\q\leq n<p}}c_nu_n\right\|
 \leq\phi(A\setminus m).
\]
Thus the ordered partial sums of the subseries form a Cauchy sequence, and the
subseries converges.

Conversely, suppose that the subseries converges and let
\[
 z_m=\sum_{\substack{n\in A\\n\geq m}}c_nu_n.
\]
A $1$-unconditional basis has the lattice-monotonicity property: if
$0\leq y\leq z$ coordinatewise, then $\|y\|\leq\|z\|$.  Hence, for every
finite $F\subseteq A\setminus m$,
\[
 \left\|\sum_{n\in F}c_nu_n\right\|\leq\|z_m\|.
\]
Therefore $\phi(A\setminus m)\leq\|z_m\|\to0$, proving
$\J_{X,c}=\Exh(\phi)$.

It remains to prove non-pathology.  First note that for finite $F$,
\begin{equation}\label{eq:finite-phi}
 \phi(F)=\left\|x_F\right\|,
 \qquad x_F=\sum_{n\in F}c_nu_n.
\end{equation}
Indeed, each positive subsum is coordinatewise dominated by $x_F$, and the
reverse inequality follows by choosing the whole set $F$.

Choose $f\in S_{X^*}$ with $f(x_F)=\|x_F\|$.  Define on the coordinate span
\[
 g(u_n)=|f(u_n)|.
\]
For a finitely supported $y=\sum_n y_nu_n$, choose unimodular scalars
$\theta_n$ so that $\theta_nf(u_n)=|f(u_n)|$.  Then
\[
 |g(y)|
 \leq\sum_n|y_n|\,|f(u_n)|
 =\left|f\left(\sum_n |y_n|\theta_nu_n\right)\right|
 \leq\left\|\sum_n |y_n|\theta_nu_n\right\|
 =\|y\|.
\]
Thus $g$ extends to a positive functional in $B_{X^*}$.  Since $x_F$ has
nonnegative coordinates,
\[
 g(x_F)=\sum_{n\in F}c_n|f(u_n)|
 \geq |f(x_F)|=\|x_F\|,
\]
and the opposite inequality follows from $\|g\|\leq1$.  Hence
$g(x_F)=\phi(F)$.

Define a finitely supported measure $\mu_F$ by
\[
 \mu_F(\{n\})=
 \begin{cases}
  c_ng(u_n),&n\in F,\\
  0,&n\notin F.
 \end{cases}
\]
For every $B\subseteq\N$,
\[
 \mu_F(B)
 =g\left(\sum_{n\in B\cap F}c_nu_n\right)
 \leq\left\|\sum_{n\in B\cap F}c_nu_n\right\|
 \leq\phi(B),
\]
while $\mu_F(F)=\phi(F)$.  Therefore, for arbitrary $A\subseteq\N$,
\[
 \sup_{\substack{\mu\text{ finitely supported}\\\mu\leq\phi}}\mu(A)
 \geq\sup_{F\in[A]^{<\omega}}\mu_F(A)
 =\sup_{F\in[A]^{<\omega}}\phi(F)
 =\phi(A).
\]
The reverse inequality is automatic from $\mu\leq\phi$.  This is exactly
\eqref{eq:nonpath-submeasure}.
\end{proof}

\begin{remark}
For a merely unconditional basis, an equivalent renorming makes it
$1$-unconditional, so the same ideal is non-pathological.
\end{remark}

\subsection{Relation to the established representability theorem}

The preceding theorem should not be presented as a new characterization.
Borodulin-Nadzieja, Farkas, and Plebanek call an ideal Banach-representable if
it has the form
\[
 \left\{A\subseteq\N:\sum_{n\in A}x_n
        \text{ converges unconditionally in a Banach space}\right\}.
\]
They proved that the Banach-representable ideals are exactly the
non-pathological analytic $P$-ideals \cite{BNFP}; the later paper
\cite{BNF} develops the connection with Banach spaces having unconditional
bases in greater depth.  Taking $x_n=c_nu_n$ in
\cref{thm:unconditional-support} gives a particularly simple positive,
coordinatewise instance of their result.

\begin{corollary}\label{cor:source-mechanism}
The positive weighted-subseries mechanism used in Theorem 7.1 of
\cite{KS} cannot realize a pathological analytic $P$-ideal merely by replacing
the Banach sequence norm while retaining an unconditional coordinate basis.
\end{corollary}

\begin{proof}
In that construction, membership is equivalent to convergence of a positive
weighted coordinate subseries.  Apply \cref{thm:unconditional-support}.
\end{proof}

Thus any realization of a pathological analytic $P$-ideal, should one exist,
must exploit genuinely different behavior: for example, non-unconditional
geometry or the finite-union structure in the definition of a critical ideal.

\section{Residual projections and the strong-operator reformulation}

Throughout the rest of the note, assume that $\CR(X,\widehat a)$ is proper and
write
\[
 \I=\CR(X,\widehat a).
\]
Biorthogonality gives a rigid algebraic structure.

\begin{lemma}[Nested projection calculus]\label{lem:projection-calculus}
For all $m,n\in\N$,
\[
 S_nS_m=S_{\min(n,m)},
 \qquad
 R_nR_m=R_{\max(n,m)}.
\]
Moreover, $X$ is separable and $\spanop\{a_n:n\in\N\}$ is dense in $X$.
\end{lemma}

\begin{proof}
The identities follow by applying the coordinate functionals to the finite
sums.  For density, fix $x\in X$ and $\eps>0$.  Since $B(x,\eps)\in\I$ and
$\I$ is proper, some $n$ lies outside $B(x,\eps)$; hence
$\|x-S_nx\|<\eps$, while $S_nx$ belongs to the finite-dimensional span of the
$a_i$.
\end{proof}

Because scaling absorbs the threshold, put
\[
 B_x=\{n:\|R_nx\|\geq1\}.
\]
Then
\begin{equation}\label{eq:critical-scaled}
 C\in\I
 \quad\Longleftrightarrow\quad
 \text{there are }F\in\Fin\text{ and }x_1,\ldots,x_r\in X
 \text{ with }C\subseteq F\cup\bigcup_{j=1}^r B_{x_j}.
\end{equation}

\begin{proposition}[Strong-operator accumulation]\label{prop:sot-accumulation}
For $A\subseteq\N$, the following are equivalent:
\begin{enumerate}[label=(\roman*)]
 \item $A\notin\I$;
 \item for every $m$, every finite family $x_1,\ldots,x_r\in X$, and every
       $\eps>0$, there is $n\in A\setminus m$ such that
       $\|R_nx_j\|<\eps$ for all $j\leq r$;
 \item $0$ is an $\omega$-accumulation point of $\{R_n:n\in A\}$ in the
       strong operator topology.
\end{enumerate}
\end{proposition}

\begin{proof}
The negation of (ii) says that, for some $m$, finite family, and $\eps>0$,
$A\setminus m$ is covered by the corresponding bad sets.  By scaling the
vectors, this is exactly membership in the ideal described by
\eqref{eq:critical-scaled}.  The equivalence of (ii) and (iii) is the
definition of strong-operator accumulation.
\end{proof}

The distinction between ordinary closure and $\omega$-accumulation is useful:
the latter records that the common good indices may always be chosen
arbitrarily far out.

\section{An exact statistical selection criterion}

For a sequence $(T_k)$ of operators, write
\[
 T_k\longrightarrow_{\Id}0\quad\text{strongly}
\]
if, for every $x\in X$ and $\eps>0$, the set
$\{k:\|T_kx\|\geq\eps\}$ has asymptotic density zero.

\begin{theorem}[Statistical SOT criterion]\label{thm:statistical-criterion}
Assume that $\I=\CR(X,\widehat a)$ is an analytic $P$-ideal.  The following are
equivalent:
\begin{enumerate}[label=(\roman*)]
 \item $\I$ is non-pathological.
 \item For every $A\notin\I$, there is a map $f:\N\to A$ such that
       \begin{enumerate}[label=(\alph*)]
        \item $f^{-1}[\{n\}]\in\Id$ for every $n\in A$;
        \item $R_{f(k)}\to_{\Id}0$ strongly.
       \end{enumerate}
 \item For every $A\notin\I$, there are finitely supported probability
       measures $(\nu_j)$ on $A$ such that
       \begin{enumerate}[label=(\alph*)]
        \item $\nu_j(\{n\})\to0$ for every fixed $n\in A$;
        \item $\nu_j(B(x,\eps))\to0$ for every $x\in X$ and $\eps>0$.
       \end{enumerate}
\end{enumerate}
\end{theorem}

\begin{proof}
Assume (i), and fix $A\notin\I$.  By \cref{thm:hrusak}, there is a map
$f:\N\to A$ witnessing $\I\restr A\leq_K\Id$.  Since every singleton belongs
to $\I\restr A$, condition (ii)(a) holds.  Every bad set
$B(x,\eps)\cap A$ belongs to $\I\restr A$, so its inverse image under $f$ has
density zero.  This is (ii)(b).

Assume (ii).  Given $A\notin\I$, define the empirical measures
\[
 \nu_j(C)=\frac{1}{j}\bigl|\{k<j:f(k)\in C\}\bigr|
 \qquad(j\geq1).
\]
They are finitely supported probabilities on $A$.  Conditions (iii)(a) and
(iii)(b) are precisely (ii)(a) and (ii)(b).

Finally, assume (iii).  For each $A\notin\I$, let
\[
 \J_A=\{C\subseteq\N:\nu_j(C)\to0\}.
\]
The rows $(\nu_j(\{n\}))_n$ form a regular nonnegative matrix: row sums are
one and each fixed column tends to zero.  Thus $\J_A$ is a regular matrix
summability ideal.  By \eqref{eq:critical-scaled}, every $C\in\I$ is contained
in a finite set together with finitely many bad sets.  Condition (iii)(a)
kills the finite set and (iii)(b) kills the bad sets, so $\I\subseteq\J_A$.
On the other hand, $\nu_j(A)=1$, so $A\notin\J_A$.  Consequently,
\[
 \I=\bigcap_{A\notin\I}\J_A.
\]
By the matrix part of \cref{thm:hrusak}, $\I\restr A\leq_K\Id$ for every
$A\notin\I$; the first part of that theorem now gives (i).
\end{proof}

\begin{remark}[What a complete negative solution would require]
By \cref{thm:statistical-criterion}, proving that every analytic $P$-ideal
critical ideal is non-pathological is equivalent to producing, on each
positive set, a diffuse statistical sampling of the residual projections that
converges strongly to zero.  This is more precise than merely observing that
$0$ is in the strong-operator closure: \cref{prop:sot-accumulation} supplies
arbitrarily good common indices for each finite family of vectors, whereas
\cref{thm:statistical-criterion} asks for one sampling that works for every
vector.
\end{remark}

\section{A bounded-positive-subset criterion}

The next criterion is elementary but useful.  It isolates one concrete place
where an ordinary uniform boundedness argument would finish the problem.

\begin{theorem}[Local boundedness implies non-pathology]
\label{thm:bounded-positive}
Assume that $\I=\CR(X,\widehat a)$ is an analytic $P$-ideal.  Suppose that for
every $A\notin\I$ there is a set $B\subseteq A$, $B\notin\I$, such that
\[
 \sup_{n\in B}\|R_n\|<\infty.
\]
Then $\I$ is non-pathological.  In fact, for every $A\notin\I$ there is
$B\subseteq A$, $B\notin\I$, such that
\[
 \I\restr B=\Fin\restr B.
\]
\end{theorem}

\begin{proof}
Fix $A\notin\I$ and such a $B$, and enumerate $B$ increasingly as
$B=\{n_k:k\in\N\}$.  Let $M=\sup_{n\in B}\|R_n\|$.  On the dense coordinate
span, $R_ny=0$ for all sufficiently large $n$.  Given $x\in X$ and $\eps>0$,
choose a finitely supported $y$ with $\|x-y\|<\eps/M$ (with the evident
modification if $M=0$).  For all sufficiently large $k$,
\[
 \|R_{n_k}x\|=\|R_{n_k}(x-y)\|\leq M\|x-y\|<\eps.
\]
Thus $R_{n_k}\to0$ strongly in the ordinary sense.

If $C\in\I$, write $C$ as in \eqref{eq:critical-scaled}.  Each bad set
$B_{x_j}$ meets $B$ only finitely because $R_{n_k}x_j\to0$.  Hence
$C\cap B$ is finite.  Therefore $\I\restr B=\Fin\restr B$.
The increasing enumeration of $B$ witnesses
$\I\restr A\leq_K\Fin\leq_K\Id$.  Apply \cref{thm:hrusak}.
\end{proof}

The hypothesis is stronger than non-pathology, but it can be verified by a
filter-dependent uniform boundedness principle.

\subsection{A rapid$^+$ corollary}

A subset of $\N$ is $\I^*$-stationary precisely when it is $\I$-positive.
A filter $\mathcal F$ is \emph{rapid$^+$} if every $\mathcal F$-stationary set
has an $\mathcal F$-stationary subset whose increasing enumeration dominates
any prescribed increasing function.  It is a \emph{weak $P^+$-filter} if,
for each decreasing sequence $(F_k)\subseteq\mathcal F$ and each stationary
$A$, there is a stationary $B\subseteq A$ with $B\subseteq^*F_k$ for every
$k$.

De Bondt and Vernaeve proved that a filter is a stationary uniform-boundedness
filter for Banach spaces exactly when it is rapid$^+$ and weak $P^+$
\cite[Theorem 5.2]{DeBondtVernaeve}.

\begin{corollary}[Rapid$^+$ dual filters]\label{cor:rapid}
Let $\I=\CR(X,\widehat a)$ be an analytic $P$-ideal.  If the dual filter
$\I^*$ is rapid$^+$, then $\I$ is non-pathological.  More strongly, every
$A\notin\I$ contains $B\notin\I$ with $\I\restr B=\Fin\restr B$.
\end{corollary}

\begin{proof}
The dual of a $P$-ideal is a weak $P^+$-filter.  Indeed, if $(F_k)$ is a
decreasing sequence in $\I^*$, the $P$-ideal property applied to their
complements gives $F\in\I^*$ with $F\subseteq^*F_k$ for all $k$; intersecting
$F$ with a stationary set preserves stationarity.

Put $\mathcal F=\I^*$.  Since $R_nx\to_{\mathcal F}0$ for every $x\in X$,
the sequence $(R_n)$ is pointwise $\mathcal F$-bounded.  By the stationary
uniform boundedness principle, for every $\mathcal F$-stationary
$A$ there is an $\mathcal F$-stationary $B\subseteq A$ on which the operators
$R_n$ are equicontinuous.  In normed spaces this means
$\sup_{n\in B}\|R_n\|<\infty$.  Apply \cref{thm:bounded-positive}.
\end{proof}

This corollary is conditional: it does not show that every dual filter arising
from an analytic $P$-ideal is rapid$^+$.  It does, however, completely settle
the source question for the subclass whose dual filters have that property.

\section{A necessary reciprocal-norm condition}

A recent theorem of Kadets and Manskova gives a different, quantitative
restriction.  In their terminology, if $(a_n)$ is an $\mathcal F$-basis with
continuous coordinate functionals in an infinite-dimensional Banach space,
then the sequence of partial-sum norms is $(\mathcal F,1)$-admissible
\cite[Theorem 4.3]{KadetsManskova}.  Translating their definition gives the
following.

\begin{proposition}[Reciprocal-norm divergence]\label{prop:reciprocal}
Assume that $X$ is infinite-dimensional and
$\I=\CR(X,\widehat a)$.  Then for every $A\notin\I$,
\[
 \sum_{n\in A}\frac{1}{\|S_n\|}=\infty.
\]
Consequently, also
\[
 \sum_{n\in A}\frac{1}{1+\|R_n\|}=\infty.
\]
In particular, the summable ideal
\[
 \I_w=\left\{A\subseteq\N:
             \sum_{n\in A}\frac{1}{1+\|R_n\|}<\infty\right\}
\]
is contained in $\I$.
\end{proposition}

\begin{proof}
Use the dual filter $\mathcal F=\I^*$.  Its stationary sets are exactly the
sets outside $\I$, so the first assertion is precisely
$(\mathcal F,1)$-admissibility.  Since $S_n$ and $R_n=I-S_n$ are nonzero
projections in an infinite-dimensional space, their norms are at least one,
and
\[
 \|R_n\|\leq1+\|S_n\|.
\]
Hence
\[
 \frac{1}{1+\|R_n\|}
 \geq \frac{1}{2+\|S_n\|}
 \geq \frac{1}{3\|S_n\|}.
\]
This proves the second divergence statement.  If $A\in\I_w$ but
$A\notin\I$, the second statement gives a contradiction; thus
$\I_w\subseteq\I$.
\end{proof}

This is a genuine restriction on every critical ideal, but it does not by
itself yield the probability measures in \cref{thm:statistical-criterion}.
Divergence of reciprocal norms allows the norms to grow without bound on every
positive set, and that is precisely the regime in which extension from a dense
set of vectors becomes difficult.

\section{The remaining selection problem}

The $P$-ideal property already gives a useful preliminary diagonalization.
Let $(x_j)$ be dense in $X$.  All sets
\[
 \{n:\|R_nx_j\|\geq 1/q\}
 \qquad(j,q\geq1)
\]
belong to $\I$.  If $A\notin\I$, the $P$-ideal property gives $C\in\I$ that
almost contains all of them.  Therefore, on the positive set $A\setminus C$,
$R_nx_j\to0$ ordinarily for every $j$.  What is missing is the passage from
the dense set $(x_j)$ to all of $X$ when $\|R_n\|$ is unbounded.

\begin{question}[Statistical strong-operator selection]
\label{q:selection}
Let $(R_n)$ be the residual projections of an ideal Schauder system, and
suppose $\I=\CR(X,\widehat a)$ is an analytic $P$-ideal.  For every
$A\notin\I$, must there exist a map $f:\N\to A$ with density-zero fibers such
that
\[
 R_{f(k)}\longrightarrow_{\Id}0
 \quad\text{strongly}?
\]
Equivalently, must there be diffuse finitely supported probabilities on $A$
for which $R_nx\to0$ in probability for every $x\in X$?
\end{question}

By \cref{thm:statistical-criterion}, an affirmative answer to
\cref{q:selection} would prove that every analytic $P$-ideal critical ideal is
non-pathological.  Since pathological analytic $P$-ideals are known to exist,
that would give a negative answer to Question 11.2 of \cite{KS}.  Conversely,
a realization of a pathological analytic $P$-ideal would necessarily furnish
a counterexample to the selection principle.

The algebra
\[
 R_nR_m=R_{\max(n,m)}
\]
is extra structure not present in an arbitrary pointwise-null sequence of
operators, and it remains a plausible source of further progress.  The
results above show exactly what that structure would have to provide: either a
bounded positive thinning, or at least a diffuse statistical thinning with
strong convergence.

\section{Literature relation and final assessment}

The relationship among the results can be summarized as follows.

\begin{enumerate}[label=\arabic*.]
 \item Kwela--Swaczyna realize every non-pathological analytic $P$-ideal as a
 critical ideal by a simple construction over an unconditional coordinate
 basis \cite[Theorem 7.1]{KS}.

 \item The packet's unconditional-support theorem is correct after the repairs
 in \cref{thm:unconditional-support}, but it is subsumed by the
 Borodulin-Nadzieja--Farkas--Plebanek characterization of Banach-representable
 ideals \cite{BNFP}.

 \item The Hru\v{s}\'ak/Filip\'ow--Tryba criterion translates the unresolved
 converse into the exact statistical SOT condition of
 \cref{thm:statistical-criterion} \cite{FilipowTryba}.

 \item De Bondt--Vernaeve's filter-dependent uniform boundedness theorem gives
 the rapid$^+$ case, \cref{cor:rapid} \cite{DeBondtVernaeve}.

 \item Kadets--Manskova's theorem gives the reciprocal-norm obstruction in
 \cref{prop:reciprocal} \cite{KadetsManskova}.

 \item Avil\'es--Rosendal--Taylor--Tradacete prove that every basis with
 respect to an analytic filter is also a basis with respect to some Borel
 filter \cite{ARTT}.  This gives a Borel envelope but does not identify the
 critical ideal and does not resolve Question 11.2.
\end{enumerate}

\begin{status}
The full realization question remains unresolved by the present argument.  The
best complete statements obtained here are:
\begin{itemize}
 \item the corrected unconditional-support obstruction;
 \item the exact statistical-SOT and matrix-averaging characterization;
 \item non-pathology under bounded positive thinning, in particular under a
 rapid$^+$ dual filter;
 \item reciprocal-norm divergence on every positive set.
\end{itemize}
No claim of novelty is made for the unconditional theorem, and the
operator-theoretic reductions should be independently checked before being
cited as new results.
\end{status}

\begin{thebibliography}{99}

\bibitem{ARTT}
A.~Avil\'es, C.~Rosendal, M.~A.~Taylor, and P.~Tradacete,
\emph{Coordinate systems in Banach spaces and lattices},
arXiv:2406.11223, version 3, 2026.

\bibitem{BNFP}
P.~Borodulin-Nadzieja, B.~Farkas, and G.~Plebanek,
\emph{Representations of ideals in Polish groups and in Banach spaces},
J. Symbolic Logic \textbf{80} (2015), no.~4, 1268--1289,
\href{https://doi.org/10.1017/jsl.2014.63}{doi:10.1017/jsl.2014.63}.


\bibitem{BNF}
P.~Borodulin-Nadzieja and B.~Farkas,
\emph{Analytic $P$-ideals and Banach spaces},
J. Funct. Anal. \textbf{279} (2020), no.~8, 108702.

\bibitem{DeBondtVernaeve}
B.~De Bondt and H.~Vernaeve,
\emph{Filter-dependent versions of the uniform boundedness principle},
J. Math. Anal. Appl. \textbf{495} (2021), no.~1, 124705,
\href{https://doi.org/10.1016/j.jmaa.2020.124705}{doi:10.1016/j.jmaa.2020.124705}.

\bibitem{FilipowTryba}
R.~Filip\'ow and J.~Tryba,
\emph{Path of pathology}, arXiv:2501.00503, 2025.

\bibitem{KadetsManskova}
V.~Kadets and M.~Manskova,
\emph{Norms of partial sums operators for a basis with respect to a filter},
J. Math. Anal. Appl. \textbf{558} (2026), no.~2, 130481,
\href{https://doi.org/10.1016/j.jmaa.2026.130481}{doi:10.1016/j.jmaa.2026.130481};
arXiv:2509.20002.

\bibitem{KS}
A.~Kwela and J.~Swaczyna,
\emph{Zoo of ideal Schauder bases},
J. Symbolic Logic (published online 2025), 1--31,
\href{https://doi.org/10.1017/jsl.2025.10136}{doi:10.1017/jsl.2025.10136}.

\bibitem{Solecki}
S.~Solecki,
\emph{Analytic ideals and their applications},
Ann. Pure Appl. Logic \textbf{99} (1999), 51--72.

\end{thebibliography}

\end{document}
